\documentclass[11pt,a4paper]{article}

\usepackage[T1]{fontenc}
\usepackage[utf8]{inputenc}
\usepackage[english]{babel}
\usepackage{amsmath,amssymb,amsthm,mathtools}
\usepackage{booktabs,array}
\usepackage[a4paper,margin=27mm]{geometry}
\usepackage[hidelinks]{hyperref}

\newtheorem{theorem}{Theorem}[section]
\newtheorem{proposition}[theorem]{Proposition}
\newtheorem{corollary}[theorem]{Corollary}
\theoremstyle{definition}
\newtheorem{definition}[theorem]{Definition}
\theoremstyle{remark}
\newtheorem{remark}[theorem]{Remark}

\newcommand{\R}{\mathbb{R}}
\newcommand{\Sact}{\mathcal{S}}
\newcommand{\Inv}{\mathcal{I}}
\newcommand{\Pair}{\mathcal{B}}

\title{Machine-Checked Variational Reconstruction and\\
Non-Uniqueness Certificates for Janus-Type Models}
\author{}
\date{Working draft -- \today}

\begin{document}
\maketitle

\begin{abstract}
We isolate and formally verify four elementary but consequential steps in the
variational analysis of Janus-type geometric models. First, fixed metric or
symplectic data on a moduli proxy do not select a potential or Hamiltonian.
Second, a two-component linear source is the Hessian of a quadratic potential
if and only if its coefficient matrix is symmetric; the remaining affine
ambiguity is removed by a critical-point condition and one normalization.
Third, exact elimination of a scalar bulk mode from a finite-rank quadratic
parent action produces a symmetric Schur-complement kernel and hence a
formally self-adjoint reduced operator. Fourth, invariant pairing families are
closed under multiplication by invariant scalar functions, so pointwise
multiplicity one does not imply a single constant global coupling. These
statements are formalized in Lean 4 and compile against Mathlib. Their purpose
is diagnostic: they separate action selection, inverse variationality,
reduction and global coupling classification. No full Janus field equation,
renormalized action or cosmological prediction is derived here.
\end{abstract}

\section{Introduction}

The passage from geometry to dynamics contains several logically distinct
problems. A metric on a parameter space can define gradients once a functional
is known, but it does not normally choose that functional. A symmetric field
equation can be compatible with a variational principle without fixing its
normalization or boundary completion. Eliminating auxiliary or bulk variables
can preserve reciprocity, yet this says nothing about whether the parent
theory was physically selected. Finally, a unique invariant tensor at one
background need not imply a unique coupling over the full background space.

These distinctions are relevant to bimetric cosmology and, in particular, to
Janus-type models in which two metric sectors and their interaction must be
specified variationally \cite{PetitDAgostini2015,PetitMargnatZejli2024}. The
present paper does not propose a replacement action. It asks a narrower
question: which finite-dimensional statements can be made exact, checked by a
small trusted kernel, and used as acceptance or rejection tests for later
field-theoretic constructions?

The inverse problem of the calculus of variations is classically described by
Helmholtz conditions and, in local field theory, by the variational bicomplex
\cite{Anderson1992,Olver1993}. Natural differential operators are likewise
organized by finite jets, equivariance and invariant theory
\cite{KolarMichorSlovak1993}. Our contribution is not a new general theorem in
those subjects. It is a machine-checked decomposition of four failure modes
that recur when one tries to infer a physical action from geometric data.

The formalization uses Lean 4 \cite{DeMouraUllrich2021} and Mathlib
\cite{Mathlib2020}. All statements reported as verified are theorem objects in
the accompanying repository. The finite models are deliberately simple: this
makes the logical conclusions transparent, while preventing them from being
mistaken for a completed continuum theory.

\section{Scope and formal method}

The results are certificates about implications between explicitly declared
data. They do not certify that the corresponding Janus manifold, bundles,
boundary conditions or quantum fields have been constructed.

We distinguish three levels:
\begin{enumerate}
  \item \emph{algebraic certificate}: a theorem proved in Lean for a finite or
        abstract model;
  \item \emph{field-theoretic interpretation}: the continuum statement that
        motivates the certificate;
  \item \emph{physical specialization}: a construction using the actual
        fields, geometry and boundary data of a Janus model.
\end{enumerate}
Only the first level is claimed in this paper. The second is discussed as an
interpretation, and the third remains future work.

The checked source is organized as four independent Lean modules. Every module
imports Mathlib, disables implicit variables, states its data structures and
proves the advertised result without axioms added by this manuscript. The
repository build command is given in Section~\ref{sec:reproducibility}.

\section{Geometry does not select a potential}
\label{sec:moduli}

Consider a one-dimensional parameter $q\in\R$ equipped with the positive
metric
\begin{equation}
  g_w = w\,dq^2, \qquad w>0.
\end{equation}
For every $a\in\R$, define
\begin{equation}
  V_a(q)=\frac{a}{2}q^2,
  \qquad
  \operatorname{grad}_{g_w}V_a(q)=\frac{a}{w}q.
\end{equation}
The metric is unchanged when $a$ changes. At $q=1$, the choices $a=1$ and
$a=2$ give distinct potential values and distinct gradients.

Likewise, let a non-degenerate constant two-form be represented by a scale
$s\ne0$. A Hamiltonian with coefficient $a$ produces the proxy slope
\begin{equation}
  X_a(q)=\frac{a}{s}q.
\end{equation}
The same symplectic scale supports distinct Hamiltonian slopes.

\begin{proposition}[Moduli-geometry non-selection]
A fixed positive metric does not determine a unique potential or gradient, and
a fixed non-degenerate symplectic scale does not determine a unique
Hamiltonian, even in the one-dimensional proxies above.
\end{proposition}

\begin{proof}
At $q=1$, one has $V_1(1)=1/2\ne1=V_2(1)$. Since $w>0$,
$1/w\ne2/w$. The symplectic statement follows similarly from $s\ne0$.
These inequalities are proved in Lean as
\path{same_metric_supports_distinct_potentials},
\path{same_metric_supports_distinct_gradients}, and
\path{same_symplectic_form_supports_distinct_hamiltonians}.
\end{proof}

The result is intentionally modest. It does not say that geometry is
irrelevant. Geometry restricts admissible functionals and turns a selected
functional into a flow. It says only that the map
\begin{equation}
  (\text{metric},\text{symplectic form},\text{complex structure})
  \longmapsto \text{one action}
\end{equation}
requires additional selection data.

\section{Quadratic inverse variational problem}
\label{sec:helmholtz}

Let a linear two-component source be represented by
\begin{equation}
  H=
  \begin{pmatrix}
    h_{xx} & h_{xy}\\
    h_{yx} & h_{yy}
  \end{pmatrix}.
\end{equation}
We ask whether $H$ is the Hessian of a quadratic potential
\begin{equation}
  V(x,y)=\frac{a}{2}x^2+bxy+\frac{c}{2}y^2+\ell_xx+\ell_yy+k.
  \label{eq:quadratic-potential}
\end{equation}

\begin{theorem}[Finite Helmholtz criterion]
There exists a potential of the form \eqref{eq:quadratic-potential} whose
Hessian is $H$ if and only if
\begin{equation}
  h_{xy}=h_{yx}.
  \label{eq:self-adjoint}
\end{equation}
\end{theorem}

\begin{proof}
Necessity follows because mixed second derivatives of
\eqref{eq:quadratic-potential} are both $b$. Conversely, if
\eqref{eq:self-adjoint} holds, choose
\begin{equation}
  a=h_{xx},\qquad b=h_{xy},\qquad c=h_{yy},
  \qquad \ell_x=\ell_y=k=0.
\end{equation}
The resulting Hessian equals $H$. Lean proves the equivalence as
\path{helmholtz_realizability_iff} and the obstruction for a non-symmetric
operator as \path{nonsymmetric_operator_is_not_variational}.
\end{proof}

The Hessian does not determine the affine part. If two quadratic potentials
have the same Hessian, their difference is
\begin{equation}
  \Delta V(x,y)=\Delta\ell_x x+\Delta\ell_y y+\Delta k.
  \label{eq:affine-ambiguity}
\end{equation}
This is checked by \path{equal_hessian_difference_affine}.

\begin{corollary}[Normalized uniqueness]
Assume \eqref{eq:self-adjoint}. Among potentials with zero linear terms and
$V(0,0)=0$, there exists exactly one potential with Hessian $H$.
\end{corollary}

\begin{proof}
Equation~\eqref{eq:affine-ambiguity} shows that the zero-linear and zero-value
conditions remove the complete affine ambiguity. The Lean theorem is
\path{pt_normalized_reconstruction_unique_existence}. In the formal module,
``PT-even'' is used only for the coefficient condition
$\ell_x=\ell_y=0$; no continuum PT transformation is inferred from it.
\end{proof}

In a field theory, formal self-adjointness is the local linear shadow of the
Helmholtz conditions. It is generally insufficient for a global action:
boundary terms, gauge identities, regularity and variational cohomology must
also be controlled \cite{Anderson1992,Olver1993}.

\section{Parent reduction and the Schur complement}
\label{sec:schur}

Let $I$ be a finite nonempty index set, let $X\in\R$ be a scalar bulk
variable, and let $q=(q_i)_{i\in I}$ be boundary variables. Consider
\begin{equation}
  \Sact(X,q)
  =\frac{a}{2}X^2+X\sum_i c_iq_i
   +\frac12\sum_{i,j}q_iB_{ij}q_j,
  \label{eq:parent-action}
\end{equation}
where $a\ne0$ and $B_{ij}=B_{ji}$.

The bulk Euler equation is
\begin{equation}
  \frac{\partial\Sact}{\partial X}=aX+\sum_i c_iq_i=0,
\end{equation}
with unique solution
\begin{equation}
  X_*(q)=-\frac{1}{a}\sum_i c_iq_i.
\end{equation}
Substitution into \eqref{eq:parent-action} yields
\begin{equation}
  \Sact_{\mathrm{red}}(q)
  =\frac12\sum_{i,j}q_i
  \left(B_{ij}-\frac{c_ic_j}{a}\right)q_j.
  \label{eq:reduced-action}
\end{equation}

\begin{theorem}[Schur-reduction certificate]
For the parent action \eqref{eq:parent-action}:
\begin{enumerate}
  \item the stationary bulk value $X_*(q)$ exists and is unique;
  \item the reduced kernel
  \begin{equation}
    H^{\mathrm{red}}_{ij}=B_{ij}-\frac{c_ic_j}{a}
  \end{equation}
  is symmetric;
  \item the associated finite-dimensional reduced operator is formally
        self-adjoint in the standard pairing;
  \item Equation~\eqref{eq:reduced-action} is the exact on-shell action.
\end{enumerate}
\end{theorem}

\begin{proof}
The first statement is division by $a\ne0$. Symmetry follows from symmetry of
$B$ and of $c_ic_j$. Reindexing the finite double sum proves formal
self-adjointness. Completing the square gives
\begin{equation}
  \Sact(X,q)-\Sact_{\mathrm{red}}(q)
  =\frac{a}{2}\bigl(X-X_*(q)\bigr)^2,
\end{equation}
which proves the on-shell identity. The corresponding Lean results are:
\begin{itemize}
  \item \path{stationary_bulk_actual_exists_unique};
  \item \path{reduced_boundary_hessian_symmetric};
  \item \path{reduced_hessian_formally_self_adjoint};
  \item \path{parent_action_on_shell};
  \item \path{parent_action_sub_reduced_action_eq_square}.
\end{itemize}
\end{proof}

For $a>0$, the square identity also makes $X_*(q)$ the unique global minimizer
at fixed $q$; for $a<0$, it is the unique global maximizer. These statements
are checked separately in the formal source. The result explains why
reciprocity survives exact quadratic elimination. It does not select the
parent coefficients $a,c_i,B_{ij}$ and does not construct a bulk PDE or its
boundary functional.

\section{Pointwise pairings versus global couplings}
\label{sec:pairings}

Let a symmetry act on a background space $\mathcal{B}$ and a field space
$E$. An invariant scalar is a function $f:\mathcal{B}\to\R$ constant on
symmetry orbits. An invariant pairing family is a map
\begin{equation}
  \beta:\mathcal{B}\times E\times E\longrightarrow\R
\end{equation}
that is unchanged when the background and both field arguments are transformed
simultaneously.

\begin{theorem}[Invariant-coefficient module]
If $f$ is an invariant scalar and $\beta$ is an invariant pairing family, then
\begin{equation}
  (f\beta)_b(u,v)=f(b)\beta_b(u,v)
\end{equation}
is an invariant pairing family.
\end{theorem}

\begin{proof}
For a symmetry element $g$, invariance gives $f(gb)=f(b)$ and
$\beta_{gb}(gu,gv)=\beta_b(u,v)$. Multiplication gives the result. This is
\path{scale_pairing_preserves_invariance} in Lean.
\end{proof}

Thus the global pairing space $\Pair$ is naturally a module over the algebra
$\Inv$ of invariant scalar functions. It need not be a one-dimensional vector
space over constant scalars.

The formal counterexample takes $\mathcal{B}=\{0,1\}$ with trivial symmetry,
$E=\R$, and
\begin{equation}
  \beta_b(u,v)=uv,
  \qquad
  f(0)=0,\quad f(1)=1,
  \qquad
  \gamma=f\beta.
\end{equation}
Both $\beta$ and $\gamma$ are invariant and have the same pointwise pairing
shape. If $\gamma=c\beta$ for one constant $c$, evaluation at $b=0$, $u=v=1$
forces $c=0$, whereas evaluation at $b=1$ forces $c=1$.

\begin{corollary}[Failure of constant global uniqueness]
Pointwise multiplicity one of the pairing shape does not imply that every
global invariant pairing is one constant multiple of a chosen pairing.
\end{corollary}

The contradiction is checked by
\path{varying_pairing_is_not_one_constant_multiple} and packaged by
\path{pointwise_shape_does_not_force_global_constant_scale}. In a geometric
application, possible invariant coefficients include scalar contractions of
curvature, second fundamental forms or gauge curvature. Their actual existence
and smooth extension across isotropy strata must be proved in the physical
category.

\section{What the certificates imply -- and what they do not}

The four results form a logical sequence, but not a derivation of a Janus
action:
\begin{equation}
\begin{aligned}
\text{geometric data}
&\not\Rightarrow \text{selected functional},\\
\text{symmetric linear source}
&\Rightarrow \text{quadratic primitive up to affine terms},\\
\text{specified quadratic parent}
&\Rightarrow \text{reciprocal reduced action},\\
\text{pointwise invariant shape}
&\not\Rightarrow \text{one constant global coupling}.
\end{aligned}
\end{equation}

For our view of the physical world, the cautious interpretation is useful.
Symmetry and geometry may strongly restrict laws without uniquely choosing
them. Apparent constants can only be promoted to universal constants after
background-dependent invariant coefficients have been excluded or fixed.
Likewise, a well-behaved reduced theory can inherit its consistency from a
parent action, while the physical origin of that parent remains an independent
question.

Nothing here establishes two metric sectors in nature, negative masses, a
specific interaction tensor, anomaly cancellation, stability of the full
field space, or an absolute cosmological scale. Those claims require concrete
geometric and analytic constructions followed by observational tests.

\section{Reproducibility}
\label{sec:reproducibility}

The formal statements are contained in the following modules:
\begin{itemize}
  \item \path{P0EFTJanusModuliGeometryNoGo.lean};
  \item \path{P0EFTJanusUniversalActionProperty.lean};
  \item \path{P0EFTJanusHessianHelmholtzReconstruction.lean};
  \item \path{P0EFTJanusFiniteRankParentSchurHelmholtz.lean};
  \item \path{P0EFTJanusInvariantCoefficientModule.lean}.
\end{itemize}
They are imported by focused Program P heads in the accompanying repository.
The principal target is
\begin{quote}\small
\path{JanusFormal.Branches.FundamentalGeometryPVariationalPrinciple}.
\end{quote}
The invariant-coefficient module is checked through the focused target
\begin{quote}\small
\path{JanusFormal.Branches.FundamentalGeometryPEInvariantPairings}.
\end{quote}
Both targets build successfully. This draft uses repository commit
\texttt{f4f5e77e}, Lean~4.31.0 and Mathlib revision \texttt{fabf563a}. The
release archive will retain the corresponding full commit identifiers.

\section{Conclusion}

We have given machine-checked finite-dimensional certificates for four
separate issues in variational reconstruction: non-selection of a functional
by moduli geometry, the symmetric-Hessian criterion and its affine ambiguity,
exact Schur reduction of a quadratic parent action, and the module structure
of global invariant pairings. The results are elementary by design. Their
scientific value lies in preventing distinct questions from being conflated
when a geometric model is promoted to a dynamical theory.

The next mathematically substantive step is not to enlarge these proxies, but
to instantiate them on a concrete Janus field bundle with an explicit local
Euler source, admissible boundary data and a proved global descent theorem.
Only then can Helmholtz, Noether, anomaly and normalization conditions test a
physical action rather than its finite-dimensional shadow.

\begin{thebibliography}{9}
\small
\setlength{\itemsep}{0.35em}
\vspace{0.5em}

\bibitem{Anderson1992}
I.~M. Anderson,
``Introduction to the variational bicomplex,''
in \emph{Mathematical Aspects of Classical Field Theory},
Contemporary Mathematics \textbf{132}, 51--73 (1992).
\href{https://doi.org/10.1090/conm/132/1188434}{doi:10.1090/conm/132/1188434}.

\bibitem{Olver1993}
P.~J. Olver,
\emph{Applications of Lie Groups to Differential Equations}, 2nd ed.,
Springer, New York (1993).

\bibitem{KolarMichorSlovak1993}
I. Kol\'a\v{r}, P.~W. Michor and J. Slov\'ak,
\emph{Natural Operations in Differential Geometry},
Springer, Berlin (1993).

\bibitem{PetitDAgostini2015}
J.-P. Petit and G. D'Agostini,
``Lagrangian derivation of the two coupled field equations in the Janus
cosmological model,''
\emph{Astrophysics and Space Science} \textbf{357}, 67 (2015).
\href{https://doi.org/10.1007/s10509-015-2250-6}{doi:10.1007/s10509-015-2250-6}.

\bibitem{PetitMargnatZejli2024}
J.-P. Petit, F. Margnat and H. Zejli,
``A bimetric cosmological model based on Andre\"i Sakharov's twin universe
approach,''
\emph{European Physical Journal C} \textbf{84}, 1226 (2024).
\href{https://doi.org/10.1140/epjc/s10052-024-13569-w}{doi:10.1140/epjc/s10052-024-13569-w}.

\bibitem{DeMouraUllrich2021}
L. de Moura and S. Ullrich,
``The Lean 4 theorem prover and programming language,''
in \emph{Automated Deduction -- CADE 28}, LNCS \textbf{12699},
625--635 (2021).
\href{https://doi.org/10.1007/978-3-030-79876-5_37}{doi:10.1007/978-3-030-79876-5\_37}.

\bibitem{Mathlib2020}
The Mathlib Community,
``The Lean mathematical library,''
in \emph{Proceedings of CPP 2020}, 367--381 (2020).
\href{https://doi.org/10.1145/3372885.3373824}{doi:10.1145/3372885.3373824}.

\end{thebibliography}

\end{document}
